\chapter{A Probabilistic View of Learning}

\section{Introduction}

In Chapter 1, we introduced learning as the problem of constructing a useful mapping from inputs to outputs based on data. That viewpoint is broad and powerful, but it leaves an important question unresolved: where do the data come from? Why should performance on a finite sample tell us anything about future performance? Why do some prediction rules make sense under uncertainty while others do not?

Probability provides the language in which these questions can be made precise. In machine learning, data are not usually treated as isolated observations. Rather, they are modeled as realizations of random variables generated by an unknown process. A learning algorithm receives a sample from this process and attempts to infer a rule that performs well on future observations drawn in a similar way. This probabilistic viewpoint allows us to define prediction error mathematically, distinguish training performance from true performance, and characterize optimal predictors under different loss functions.

This chapter develops the statistical learning setup carefully. We begin by treating data as samples from an unknown distribution on input--output pairs. We then define loss functions, expected risk, and empirical risk. These concepts formalize the idea that learning is not just about fitting past data, but about minimizing expected future error. We next study conditional distributions and derive the form of optimal predictors for common losses. This leads naturally to the notions of conditional expectation, conditional median, and Bayes decision rules. We also explain the role of noise, ambiguity, and irreducible uncertainty, which show why perfect prediction is often impossible even in principle.

Finally, we discuss two important interpretive frameworks: the frequentist and Bayesian viewpoints. These are not merely philosophical labels. They reflect different ways of representing uncertainty and different ways of understanding what a model parameter means. Both perspectives play an important role in modern machine learning.

The main goal of this chapter is to move from the informal language of ``predicting from data'' to a mathematically precise statistical framework. This framework will serve as the foundation for nearly everything that follows: generalization theory, empirical risk minimization, probabilistic modeling, neural network training, and generative modeling.

\section{Data as Samples from an Unknown Distribution}

A central assumption in statistical learning is that the data we observe are generated according to some underlying probability distribution. This assumption does not mean that the world is simple or that the distribution is known. On the contrary, the distribution is usually highly complex and entirely unknown. The point is that probability gives us a way to model variability, uncertainty, and repeated sampling.

\subsection{Random Variables and Observations}

Let $X$ be a random variable taking values in an input space $\mathcal{X}$, and let $Y$ be a random variable taking values in an output space $\mathcal{Y}$. The pair $(X,Y)$ is assumed to have some joint distribution $P$ on $\mathcal{X}\times\mathcal{Y}$.

A dataset of size $n$ is modeled as a sample
\[
(X_1,Y_1), (X_2,Y_2), \dots, (X_n,Y_n),
\]
typically assumed to be independent and identically distributed (i.i.d.) according to $P$:
\[
(X_i,Y_i) \sim P, \qquad i=1,\dots,n.
\]
The i.i.d.\ assumption is not universally valid, but it is the basic starting point for much of statistical learning theory.

\begin{definition}[Statistical Learning Setup]
A statistical learning problem consists of a pair of random variables $(X,Y)$ with unknown joint distribution $P$ on $\mathcal{X}\times\mathcal{Y}$, together with a hypothesis class $\mathcal{H}$ of functions $f:\mathcal{X}\to\mathcal{A}$, where $\mathcal{A}$ is the prediction space, and a loss function
\[
\ell:\mathcal{A}\times\mathcal{Y}\to[0,\infty).
\]
The goal is to choose a predictor $f\in\mathcal{H}$ that performs well under the distribution $P$.
\end{definition}

Here $\mathcal{A}$ is the set of possible predictions. In regression, one often has $\mathcal{A}=\mathbb{R}$. In classification, $\mathcal{A}$ may be a finite label set or a probability simplex over labels.

\subsection{Why Distributions Matter}

The probabilistic setup allows us to formalize a crucial idea: the learner is not judged only on the finite training sample, but on future examples generated from the same underlying process. That is why the distribution $P$ matters. It determines what kinds of inputs are likely, how outputs relate to inputs, and what uncertainty is intrinsic to the problem.

A single finite sample provides only partial information about $P$. Learning, therefore, is the problem of drawing reliable conclusions about future performance from incomplete evidence.

\subsection{Marginals, Conditionals, and the Joint Law}

The joint distribution $P$ can be understood through:
\begin{itemize}
    \item the marginal distribution of $X$,
    \item the marginal distribution of $Y$,
    \item the conditional distribution of $Y$ given $X$.
\end{itemize}

In supervised learning, the conditional law of $Y$ given $X$ is especially important. It describes how the output behaves for a given input. In many tasks, the ideal predictor is determined entirely by this conditional distribution.

For example:
\begin{itemize}
    \item in regression with squared loss, the optimal predictor is the conditional mean,
    \item in regression with absolute loss, the optimal predictor is the conditional median,
    \item in classification with zero-one loss, the optimal predictor is the most probable class.
\end{itemize}

These facts reveal that prediction is fundamentally about conditional probability.

\begin{example}[Medical Prediction]
Suppose $X$ represents patient features and $Y$ represents a diagnosis or outcome. Then the learning problem is governed not only by the overall frequency of the disease, but by the conditional distribution of $Y$ given $X$. Two patients may have different risks because their feature vectors correspond to different conditional distributions.
\end{example}

\section{Loss Functions and Risk}

To judge the quality of a prediction, we need a quantitative notion of error. This is supplied by a loss function.

\begin{definition}[Loss Function]
A loss function is a map
\[
\ell:\mathcal{A}\times\mathcal{Y}\to[0,\infty)
\]
such that $\ell(a,y)$ measures the penalty incurred by predicting $a$ when the true output is $y$.
\end{definition}

The choice of loss depends on the task and on what kinds of mistakes matter.

\subsection{Common Loss Functions}

Some standard examples are:
\begin{itemize}
    \item \textbf{Squared loss:}
    \[
    \ell(a,y) = (a-y)^2.
    \]
    \item \textbf{Absolute loss:}
    \[
    \ell(a,y) = |a-y|.
    \]
    \item \textbf{Zero-one loss for classification:}
    \[
    \ell(a,y)=
    \begin{cases}
        0, & a=y,\\
        1, & a\neq y.
    \end{cases}
    \]
    \item \textbf{Log loss / cross-entropy:} if $a$ is a probability distribution over classes and $y$ is the true class, then
    \[
    \ell(a,y) = -\log a(y).
    \]
\end{itemize}

Different loss functions lead to different notions of optimality. A predictor that is optimal under squared loss need not be optimal under absolute loss or zero-one loss.

\subsection{Population Risk}

The most important quantity in statistical learning is expected loss under the data-generating distribution.

\begin{definition}[Population Risk]
Let $(X,Y)\sim P$, and let $f:\mathcal{X}\to\mathcal{A}$ be a predictor. The population risk of $f$ is
\[
R(f) = \mathbb{E}[\ell(f(X),Y)].
\]
\end{definition}

This quantity represents the average loss the predictor would incur on future data drawn from the same distribution.

Population risk is the true performance criterion. It is what we actually care about. However, it is usually not directly computable because the distribution $P$ is unknown.

\subsection{Empirical Risk}

Given a sample $(X_i,Y_i)_{i=1}^n$, the empirical risk of $f$ is
\[
\widehat{R}_n(f) = \frac{1}{n}\sum_{i=1}^n \ell(f(X_i),Y_i).
\]
Unlike population risk, this quantity is observable. Learning algorithms often minimize empirical risk or a regularized version of it.

The distinction between these two risks is fundamental:
\begin{itemize}
    \item \emph{population risk} measures true expected future performance;
    \item \emph{empirical risk} measures performance on the observed sample.
\end{itemize}

One of the main questions of learning theory is when small empirical risk implies small population risk.

\subsection{Risk Minimization}

In an ideal world, if we knew the distribution $P$, we would seek a predictor minimizing population risk:
\[
f^\star \in \arg\min_f R(f).
\]
Such a predictor is often called a \emph{Bayes predictor} or \emph{Bayes decision rule}, depending on the context.

In practice, we usually restrict attention to a hypothesis class $\mathcal{H}$ and instead try to solve
\[
\min_{f\in\mathcal{H}} R(f),
\]
or more realistically, we minimize empirical risk:
\[
\min_{f\in\mathcal{H}} \widehat{R}_n(f).
\]

Thus statistical learning can be viewed as the attempt to approximate the population risk minimizer using only a finite sample.

\section{Conditional Distributions and Optimal Prediction}

The conditional distribution of $Y$ given $X=x$ encodes the predictive structure of the problem. It tells us, for each input $x$, how the output behaves probabilistically. Many optimal prediction rules can be derived by minimizing conditional expected loss pointwise in $x$.

\subsection{Conditional Risk}

Let $f$ be a predictor. Then its population risk can be decomposed as
\[
R(f)=\mathbb{E}\big[\ell(f(X),Y)\big]
= \mathbb{E}\Big[\mathbb{E}\big[\ell(f(X),Y)\mid X\big]\Big].
\]
This is an instance of the tower property of conditional expectation.

For a fixed input value $x$, define the conditional risk of predicting $a\in\mathcal{A}$ by
\[
r_x(a) = \mathbb{E}[\ell(a,Y)\mid X=x].
\]
Then the overall risk of a predictor $f$ is
\[
R(f) = \mathbb{E}[r_X(f(X))].
\]

This leads to a fundamental principle:

\begin{proposition}[Pointwise Risk Minimization]
Suppose for each $x\in\mathcal{X}$ we choose a prediction $f^\star(x)$ satisfying
\[
f^\star(x)\in \arg\min_{a\in\mathcal{A}} \mathbb{E}[\ell(a,Y)\mid X=x].
\]
Then $f^\star$ minimizes the population risk over all measurable predictors.
\end{proposition}

\begin{proof}[Proof sketch]
For any predictor $f$, one has
\[
\mathbb{E}[\ell(f(X),Y)\mid X=x] \ge \inf_{a\in\mathcal{A}} \mathbb{E}[\ell(a,Y)\mid X=x]
\]
for each $x$. Integrating over the distribution of $X$ yields
\[
R(f)\ge R(f^\star).
\]
Thus minimizing conditional risk pointwise also minimizes total population risk.
\end{proof}

This proposition is conceptually central. It shows that the Bayes-optimal predictor is determined by the conditional distribution of $Y$ given $X$.

\section{Conditional Expectation and Optimal Prediction Under Squared Loss}

We now derive one of the most important results in statistics and machine learning: under squared loss, the optimal predictor is the conditional mean.

\begin{theorem}[Bayes Optimal Predictor Under Squared Loss]
Suppose $\mathcal{Y}\subseteq \mathbb{R}$ and the loss is
\[
\ell(a,y)=(a-y)^2.
\]
Then the predictor that minimizes population risk over all measurable functions is
\[
f^\star(x)=\mathbb{E}[Y\mid X=x].
\]
\end{theorem}

\begin{proof}
Fix $x$ and consider the conditional risk
\[
r_x(a)=\mathbb{E}\big[(a-Y)^2\mid X=x\big].
\]
Let
\[
m(x)=\mathbb{E}[Y\mid X=x].
\]
We expand:
\[
(a-Y)^2 = (a-m(x)+m(x)-Y)^2.
\]
Thus
\[
(a-Y)^2 = (a-m(x))^2 + 2(a-m(x))(m(x)-Y) + (m(x)-Y)^2.
\]
Taking conditional expectation given $X=x$ yields
\[
r_x(a)
= (a-m(x))^2 + 2(a-m(x))\mathbb{E}[m(x)-Y\mid X=x]
+ \mathbb{E}[(m(x)-Y)^2\mid X=x].
\]
Since $m(x)=\mathbb{E}[Y\mid X=x]$, we have
\[
\mathbb{E}[m(x)-Y\mid X=x]=m(x)-\mathbb{E}[Y\mid X=x]=0.
\]
Therefore
\[
r_x(a)
= (a-m(x))^2 + \mathbb{E}[(m(x)-Y)^2\mid X=x].
\]
The second term does not depend on $a$, so $r_x(a)$ is minimized when
\[
a=m(x)=\mathbb{E}[Y\mid X=x].
\]
By pointwise risk minimization, the optimal predictor is
\[
f^\star(x)=\mathbb{E}[Y\mid X=x].
\]
\end{proof}

\subsection{Interpretation}

This theorem explains why regression models trained with squared loss aim to estimate a conditional mean. Even when the data are noisy, the best squared-error predictor is not the noiseless underlying mechanism in any absolute sense; it is the average response conditioned on the input.

This also clarifies an important modeling issue. If the conditional distribution of $Y$ given $X=x$ is multimodal, then the conditional mean may not be a typical outcome. For example, if a text prompt admits several very different plausible images, the mean of all such images may look blurry and unrealistic. This is one reason why mean prediction is not always a good generative strategy, even though it is optimal under squared loss.

\subsection{Bias--Variance Style Decomposition}

The proof above also gives a decomposition:
\[
\mathbb{E}[(a-Y)^2\mid X=x]
= (a-\mathbb{E}[Y\mid X=x])^2 + \operatorname{Var}(Y\mid X=x).
\]
This shows that the conditional variance is the unavoidable residual error at input $x$. Even the Bayes-optimal predictor cannot reduce this part.

\section{Conditional Median and Optimal Prediction Under Absolute Loss}

The optimal predictor changes if we change the loss function. Under absolute loss, the correct target is no longer the conditional mean but the conditional median.

\begin{definition}[Conditional Median]
A value $m(x)$ is called a conditional median of $Y$ given $X=x$ if
\[
\mathbb{P}(Y\le m(x)\mid X=x)\ge \frac12
\qquad\text{and}\qquad
\mathbb{P}(Y\ge m(x)\mid X=x)\ge \frac12.
\]
\end{definition}

\begin{theorem}[Bayes Optimal Predictor Under Absolute Loss]
Suppose $\mathcal{Y}\subseteq\mathbb{R}$ and the loss is
\[
\ell(a,y)=|a-y|.
\]
Then any conditional median of $Y$ given $X=x$ minimizes the conditional expected loss
\[
\mathbb{E}[|a-Y|\mid X=x].
\]
Hence a Bayes-optimal predictor under absolute loss is any conditional median.
\end{theorem}

\begin{proof}[Proof sketch]
Fix $x$ and define
\[
r_x(a)=\mathbb{E}[|a-Y|\mid X=x].
\]
The function $r_x(a)$ is convex in $a$. When differentiable, its derivative is
\[
\frac{d}{da} r_x(a)
= \mathbb{P}(Y\le a\mid X=x) - \mathbb{P}(Y\ge a\mid X=x),
\]
up to technicalities at atoms. A minimizer therefore occurs when the probability mass on either side of $a$ is at least one-half, which is exactly the definition of a median. Thus any conditional median minimizes conditional expected absolute deviation.
\end{proof}

\subsection{Interpretation}

The conditional median is more robust to outliers than the conditional mean. Squared loss penalizes large deviations heavily, so it is sensitive to rare extreme values. Absolute loss grows only linearly, so it is less strongly affected by outliers.

This is a recurring theme: the loss function determines what feature of the conditional distribution is being estimated.

\begin{example}[Mean vs Median]
Suppose that given a certain input $x$, the output $Y$ equals $0$ with probability $0.9$ and equals $100$ with probability $0.1$. Then
\[
\mathbb{E}[Y\mid X=x] = 10,
\]
while the conditional median is $0$. The mean is optimal under squared loss, but the median is optimal under absolute loss.
\end{example}

\section{Bayes Classification Under Zero-One Loss}

Classification provides another important example of Bayes-optimal prediction. Here the output space is finite, say
\[
\mathcal{Y} = \{1,2,\dots,K\},
\]
and the loss is zero-one loss:
\[
\ell(a,y)=
\begin{cases}
0, & a=y,\\
1, & a\neq y.
\end{cases}
\]

\begin{theorem}[Bayes Classifier Under Zero-One Loss]
Under zero-one loss, a predictor $f^\star$ minimizing population risk is given by
\[
f^\star(x)\in \arg\max_{k\in\{1,\dots,K\}} \mathbb{P}(Y=k\mid X=x).
\]
That is, the Bayes classifier chooses the most probable class given the input.
\end{theorem}

\begin{proof}
Fix $x$ and consider predicting class $a$. The conditional risk is
\[
r_x(a)=\mathbb{E}[\ell(a,Y)\mid X=x].
\]
Since zero-one loss is $1$ unless $a=Y$, we have
\[
r_x(a)=1-\mathbb{P}(Y=a\mid X=x).
\]
Thus minimizing $r_x(a)$ is equivalent to maximizing $\mathbb{P}(Y=a\mid X=x)$. Therefore any most probable class is conditionally optimal. By pointwise risk minimization, this yields a population risk minimizer.
\end{proof}

\subsection{Bayes Error Rate}

Even the Bayes classifier may make errors, because the label may be intrinsically uncertain given the input. The minimum achievable classification error is called the \emph{Bayes error rate}. At input $x$, the conditional Bayes error is
\[
1 - \max_k \mathbb{P}(Y=k\mid X=x).
\]
Averaging over $X$ yields the total Bayes error rate.

This quantity represents irreducible uncertainty in the classification problem. If two classes overlap in feature space, then no classifier can separate them perfectly.

\section{Noise, Uncertainty, and Irreducible Error}

One of the most important lessons of the probabilistic viewpoint is that imperfect prediction may be unavoidable even with infinite data and unlimited computation.

\subsection{Aleatoric Uncertainty}

Some uncertainty is intrinsic to the data-generating process. This is often called \emph{aleatoric uncertainty}. It reflects randomness, ambiguity, or unobserved factors. For example:
\begin{itemize}
    \item two patients with the same measured features may have different outcomes,
    \item stock prices may fluctuate due to unpredictable events,
    \item multiple translations of a sentence may all be valid.
\end{itemize}

Aleatoric uncertainty cannot be eliminated by collecting more data about the same variables alone. It is built into the conditional distribution of $Y$ given $X$.

\subsection{Irreducible Error in Regression}

Under squared loss, the Bayes-optimal predictor is the conditional mean, and the minimum conditional risk is
\[
\mathbb{E}\Big[\big(Y-\mathbb{E}[Y\mid X]\big)^2 \,\Big|\, X\Big]
= \operatorname{Var}(Y\mid X).
\]
Averaging over $X$, the minimum achievable population risk is
\[
\mathbb{E}[\operatorname{Var}(Y\mid X)].
\]
This term represents the irreducible error.

\subsection{Irreducible Error in Classification}

Under zero-one loss, the irreducible error is the Bayes error rate:
\[
R^\star = \mathbb{E}\Big[1-\max_k \mathbb{P}(Y=k\mid X)\Big].
\]
Again, this is the error of the best possible classifier, not just the best classifier in a restricted hypothesis class.

\subsection{Epistemic Uncertainty}

In addition to aleatoric uncertainty, there is uncertainty due to limited knowledge. This is often called \emph{epistemic uncertainty}. It arises because the learner does not know the true distribution $P$ and only sees finitely many examples. Unlike aleatoric uncertainty, epistemic uncertainty can in principle be reduced by more data, better models, or stronger prior knowledge.

This distinction is useful:
\begin{itemize}
    \item aleatoric uncertainty is inherent in the problem;
    \item epistemic uncertainty reflects incomplete knowledge of the problem.
\end{itemize}

\begin{example}[Noisy Regression Model]
Suppose
\[
Y = g(X) + \varepsilon,
\]
where $\varepsilon$ is random noise with
\[
\mathbb{E}[\varepsilon\mid X]=0.
\]
Then
\[
\mathbb{E}[Y\mid X]=g(X),
\]
so the conditional mean recovers the signal. But unless $\varepsilon=0$ almost surely, even the best predictor incurs nonzero error because of the noise term.
\end{example}

\section{Deterministic and Probabilistic Prediction Rules}

The results above show that optimal point predictions depend on the choice of loss. But there are situations where a single deterministic prediction is not enough. In such cases, it is often more informative to predict a full conditional distribution.

\subsection{Point Prediction}

A deterministic predictor outputs a single value
\[
f(x)\in\mathcal{A}.
\]
Examples:
\begin{itemize}
    \item a predicted house price,
    \item a single label,
    \item a single next token chosen greedily.
\end{itemize}

Point prediction is appropriate when the task requires a single action or estimate, or when the uncertainty is mild.

\subsection{Probabilistic Prediction}

A probabilistic predictor outputs a distribution
\[
q(\cdot \mid x)
\]
over possible outputs. This can express uncertainty, multimodality, and confidence.

Examples:
\begin{itemize}
    \item class probabilities in classification,
    \item predictive intervals in regression,
    \item next-token distributions in language modeling,
    \item distributions over future trajectories in forecasting.
\end{itemize}

Probabilistic prediction is often strictly more informative than deterministic prediction. A deterministic rule can be obtained from a probabilistic predictor by taking a mean, median, mode, or decision that minimizes expected loss.

\subsection{Why Probabilities Matter}

A model that outputs only a class label may hide the difference between:
\begin{itemize}
    \item a case where one class has probability $0.99$,
    \item a case where it has probability $0.51$.
\end{itemize}
The predicted label is the same, but the uncertainty is very different.

In high-stakes domains such as medicine, finance, and scientific prediction, this distinction is often crucial.

\section{Bayesian and Frequentist Interpretations}

The probabilistic language of learning admits multiple interpretations. Two of the most important are the frequentist and Bayesian viewpoints.

\subsection{Frequentist View}

In the frequentist interpretation:
\begin{itemize}
    \item the data are random,
    \item the parameters of the model are fixed but unknown,
    \item probability describes long-run variability of sampling.
\end{itemize}

For example, in linear regression one may write
\[
Y = X^\top \theta + \varepsilon,
\]
where $\theta$ is an unknown fixed parameter vector. Estimation seeks to infer $\theta$ from random data.

From this perspective, uncertainty statements concern the behavior of estimators across repeated samples. Confidence intervals and sampling distributions are frequentist tools.

\subsection{Bayesian View}

In the Bayesian interpretation:
\begin{itemize}
    \item unknown quantities such as parameters are themselves treated probabilistically,
    \item prior beliefs are encoded in a prior distribution,
    \item observed data update the prior to a posterior distribution.
\end{itemize}

If $\theta$ denotes a model parameter, Bayes' rule gives
\[
p(\theta\mid \mathcal{D}) \propto p(\mathcal{D}\mid \theta)\,p(\theta),
\]
where $\mathcal{D}$ is the dataset.

From the posterior one can derive predictive distributions:
\[
p(y\mid x,\mathcal{D}) = \int p(y\mid x,\theta)\,p(\theta\mid \mathcal{D})\,d\theta.
\]

This approach naturally quantifies uncertainty due to limited data and allows prior knowledge to be incorporated in a principled way.

\subsection{Comparing the Two Views}

The frequentist and Bayesian viewpoints often lead to similar practical procedures, but they differ conceptually.

\paragraph{Frequentist emphasis.}
How does an estimator behave across repeated sampling from the population?

\paragraph{Bayesian emphasis.}
Given the observed data, what should we believe about unknown quantities?

Neither framework exhausts all of machine learning, but both are fundamental. Classical statistics developed many frequentist tools, while modern probabilistic machine learning often incorporates Bayesian ideas, especially in uncertainty quantification, latent-variable modeling, and probabilistic inference.

\subsection{A Common Ground}

Despite their differences, both perspectives share the same basic statistical learning setup:
\begin{itemize}
    \item data arise from a probabilistic mechanism,
    \item prediction quality is judged by expected loss,
    \item future performance depends on the underlying data-generating process.
\end{itemize}

Thus the probabilistic view of learning is broader than the Bayesian/frequentist divide. It is the common framework within which both can be understood.

\section{Worked Derivations and Illustrative Examples}

\subsection{Why the Conditional Mean Minimizes Squared Loss}

Let $a$ be a candidate prediction at input $x$. Then
\[
\mathbb{E}[(Y-a)^2\mid X=x]
\]
is minimized at
\[
a=\mathbb{E}[Y\mid X=x].
\]
The proof was given earlier, but its meaning is worth repeating: squared loss asks the predictor to aim for the center of mass of the conditional distribution.

\subsection{Why the Conditional Median Minimizes Absolute Loss}

Similarly,
\[
\mathbb{E}[|Y-a|\mid X=x]
\]
is minimized at a conditional median. Absolute loss asks the predictor to balance the probability mass on either side of the prediction.

\subsection{Classification as Choosing the Most Likely Class}

Under zero-one loss, the expected conditional error of predicting class $k$ is
\[
1-\mathbb{P}(Y=k\mid X=x),
\]
so the best class is simply the one with highest conditional probability.

\begin{example}[Three Different Optimal Predictors]
Suppose that for a fixed input $x$, the conditional distribution of $Y$ is:
\[
Y=
\begin{cases}
0, & \text{with probability } 0.4,\\
1, & \text{with probability } 0.2,\\
10, & \text{with probability } 0.4.
\end{cases}
\]
Then:
\begin{itemize}
    \item the conditional mean is
    \[
    \mathbb{E}[Y\mid X=x] = 0(0.4)+1(0.2)+10(0.4)=4.2,
    \]
    so this is optimal under squared loss;
    \item a conditional median can be any value in $[1,10]$ only if the cumulative distribution permits it, but here the smallest median is $1$ and the set of medians collapses to values where cumulative mass reaches one-half; in this discrete case, $1$ is a median;
    \item if the outputs were treated as class labels, then the most probable class would be tied between $0$ and $10$.
\end{itemize}
Thus the optimal prediction depends fundamentally on the loss and on how the problem is formulated.
\end{example}

\section{A Unifying Perspective}

The probabilistic view of learning adds a crucial layer to the function-approximation viewpoint developed in Chapter 1. Instead of thinking of data as arbitrary examples, we model them as realizations from an unknown distribution. Instead of measuring performance only on a sample, we evaluate predictors by expected loss under that distribution. Instead of speaking vaguely about ``best prediction,'' we derive optimal rules from conditional distributions and loss functions.

This framework leads to several powerful ideas:
\begin{itemize}
    \item learning is about minimizing expected risk, not just fitting observed data;
    \item optimal prediction depends on the loss function;
    \item conditional expectation, conditional median, and class probabilities are not arbitrary quantities but Bayes-optimal targets under different losses;
    \item some uncertainty is irreducible, because the output may remain random even after conditioning on the input;
    \item probabilistic modeling provides the bridge between observed data and principled prediction.
\end{itemize}

A useful slogan is:
\[
\text{prediction} = \text{decision under uncertainty}.
\]
The role of probability is to describe that uncertainty; the role of the loss function is to specify what counts as a good decision.

\section{Summary}

In this chapter we introduced the statistical learning framework. We modeled data as samples from an unknown joint distribution of random variables $(X,Y)$. We defined loss functions, population risk, and empirical risk, and emphasized that population risk is the true target of learning.

We then showed how the conditional distribution of $Y$ given $X$ determines optimal prediction. Under squared loss, the Bayes-optimal predictor is the conditional mean. Under absolute loss, it is a conditional median. Under zero-one loss, the Bayes classifier chooses the most probable class. These results illustrate that ``optimal prediction'' depends on the chosen loss and on the underlying uncertainty structure.

We also discussed noise and irreducible error, distinguishing uncertainty inherent in the problem from uncertainty due to limited knowledge. Finally, we compared Bayesian and frequentist interpretations of the probabilistic setup.

This chapter provides the probabilistic backbone for the rest of the book. The next chapter will examine hypothesis spaces and inductive bias, explaining why data alone do not determine a predictor and why structural assumptions are essential for generalization.

\section{Exercises}

\begin{enumerate}
    \item Let $(X,Y)$ be a random pair with loss function $\ell$. Write the definition of population risk and explain in words what it means.

    \item Let $(X_i,Y_i)_{i=1}^n$ be a dataset. Write the empirical risk of a predictor $f$ and compare it with the population risk.

    \item Suppose $Y\in\mathbb{R}$ and the loss is squared loss. Show that for a fixed random variable $Y$, the quantity
    \[
    \mathbb{E}[(Y-a)^2]
    \]
    is minimized over $a\in\mathbb{R}$ at $a=\mathbb{E}[Y]$.

    \item Extend the previous exercise to the conditional setting and show that
    \[
    f^\star(x)=\mathbb{E}[Y\mid X=x]
    \]
    minimizes the conditional expected squared loss.

    \item Suppose the loss is absolute loss. Prove, or give a careful argument, that a median minimizes expected absolute deviation.

    \item For each of the following losses, identify the Bayes-optimal predictor:
    \begin{enumerate}
        \item squared loss,
        \item absolute loss,
        \item zero-one loss.
    \end{enumerate}

    \item Let $\mathcal{Y}=\{1,2,3\}$. Suppose that for some fixed $x$,
    \[
    \mathbb{P}(Y=1\mid X=x)=0.2,\qquad
    \mathbb{P}(Y=2\mid X=x)=0.5,\qquad
    \mathbb{P}(Y=3\mid X=x)=0.3.
    \]
    What is the Bayes classifier at $x$ under zero-one loss? What is the conditional Bayes error at $x$?

    \item Give an example where the conditional mean and conditional median are different. Explain why this happens.

    \item Suppose
    \[
    Y = g(X) + \varepsilon
    \]
    with $\mathbb{E}[\varepsilon\mid X]=0$. Show that
    \[
    \mathbb{E}[Y\mid X]=g(X).
    \]

    \item Show that under squared loss, the minimum achievable conditional risk at input $x$ is
    \[
    \operatorname{Var}(Y\mid X=x).
    \]

    \item Explain the difference between aleatoric and epistemic uncertainty. Give one example of each.

    \item Consider a classifier that outputs only the predicted class label and another classifier that outputs a full vector of class probabilities. Why might the second be more informative?

    \item In your own words, explain why the optimal predictor depends on the loss function.

    \item Let $Y$ be a Bernoulli random variable with
    \[
    \mathbb{P}(Y=1\mid X=x)=p(x).
    \]
    Under zero-one loss, what is the optimal classifier? Under squared loss for predicting $Y\in\{0,1\}$ numerically, what is the optimal predictor?

    \item Compare the frequentist and Bayesian interpretations of the statement ``the parameter $\theta$ is uncertain.'' What does this mean in each framework?

    \item A model predicts the same class label for two different inputs, but in one case its estimated probability is $0.99$ and in the other case $0.55$. Why might these two predictions be treated differently in practice?

    \item Suppose that for a fixed input $x$, the conditional distribution of $Y$ is symmetric around a value $m$. Show that $m$ is both the conditional mean and a conditional median.

    \item Let $f^\star$ be a Bayes-optimal predictor over all measurable functions, and let $\mathcal{H}$ be a restricted hypothesis class. Explain why the best predictor in $\mathcal{H}$ may still have larger risk than $f^\star$.
\end{enumerate}

\section*{Historical Note}

The probabilistic formulation of learning draws from classical statistics, decision theory, and probability. The idea of evaluating predictors by expected loss is central to statistical decision theory, while concepts such as conditional expectation and Bayes rules are foundational in probability and mathematical statistics. Modern machine learning inherits these ideas and extends them to high-dimensional models, large-scale optimization, and complex predictive tasks. Although the models have changed dramatically, the underlying probabilistic logic remains one of the deepest unifying principles of the field.